\section{Conclusion}
\label{sec:conclusion}
Ce stage a eu pour objectif de finir le développement du parseur et la sérialisation vers le format XMI, d'implémenter des outils de la théorie du contrôle, des filtres et des anti-windups pour CHIPS, ainsi que d'unifier tout cela pour avoir une chaîne d'outils complète et fonctionnelle.

\subsection{Bilan technique}
\label{sec:bilan_technique}

D'un point de vue technique, le parseur est maintenant fonctionnel et on peut générer des fichiers XMI, conforme au méta-modèle, à partir du programme CHIPS. \\

Au total, 19 classes C++ en rapport avec la théorie du contrôle ont été implémentées, que ce soit des filtres ou des anti-windups. 10 bancs de tests ont été réalisés, avec pour chacun d'eux des graphiques et des résultats de simulation. L'ensemble de ces outils peuvent être utilisés avec CHIPS car nous avons mis en place une chaîne d'outils complète et fonctionnelle, qui permet de générer un programme BIP à partir d'un programme CHIPS, puis de générer du code C++ à partir de ce programme BIP, et enfin d'exécuter le programme pour obtenir les résultats de simulation. \\

Tout cela correspond à un volume de d'un peu plus de 10000 lignes de code C++, python et script bash. \\

\subsection{Les perspectives}
\label{sec:perspectives}

Différentes perspectives d'amélioration sont possibles pour ce projet. Tout d'abord, il serait intéressant de mettre en place des tests de non-régression pour que les modifications futures du parseur ne cassent pas les fonctionnalités existantes. \\

Ensuite pour la partie des filtres et des anti-windups, il faudrait utiliser le même type de données pour les variables pour éviter les différences de résultats engendrés par les différences de précision. Sur la même catégorie, il faudrait également implémenter des tests unitaires pour les classes de la théorie du contrôle pour s'assurer de leur bon fonctionnement. \\

Aussi, de pouvoir utiliser les filtres et les anti-windups directement dans un programme CHIPS, tout cela a déjà été réfléchi mais par manque de temps, cela n'a pas pu être implémenté. \\

Enfin, il serait intéressant de pouvoir avoir d'autres outils de la théorie du contrôle, comme le gain scheduling, le contrôle par avance/retard de phase\dots

% lead lag : Livre de référence: Nise, Norman S. (2004); Control Systems Engineering (4 ed.); Wiley & Sons; ISBN 0-471-44577-0
%des slides explicatifs: https://homepages.laas.fr/fgouaisb/donnees/M1ICM/slidesM1ICMp8.pdf 
% gain scheduling : Une revue de l'état de l'art: Rugh, W. J., & Shamma, J. S. (2000). Research on gain scheduling. Automatica, 36(10), 1401-1425

\newpage

\subsection{Bilan personnel}
\label{sec:bilan_personnel}

D'un point de vue personnel, ce stage s'est inscrit dans la suite logique de mon projet semestriel, débuté depuis 
octobre 2025, qui avait été l'implémentation du parseur et la sérialisation vers XMI avec deux autres personnes. 
Cela a été pour moi le plus long projet sur lequel j'ai pu travailler car il a duré au total 9 mois. \\

Ce stage m'a permis de développer mes compétences techniques en C++, ainsi que sur la modélisation d'un langage de 
programmation et de comprendre tous les tenants et aboutissant. \\

Cette expérience m'a permis de me rendre compte de l'importance de la communication lors des comptes rendus 
d'avancement, car cela permet de faire le point sur l'avancement du projet et de discuter des problèmes 
rencontrés. \\

Cela m'a également permis de mettre en pratique ainsi qu'approfondir mes connaissances sur le cours de Systèmes Cyber-Physiques (M2). \\ 

\FloatBarrier

\newpage